% !TEX root = ../main.tex
\chapter{Introduction}
\label{ch:introduction}

\section{Motivation}

Sparse matrix--dense matrix multiplication (SpMM) is one of the fundamental
computational kernels in modern machine learning and graph-based computing.
In graph neural networks (GNNs) it embodies the neighbourhood aggregation
step, in which a sparse adjacency matrix is multiplied against a dense
feature matrix~\citep{graphsage2017, gat2018}. In large-scale recommender
systems, collaborative filtering relies on sparse user--item interaction
matrices whose dimensions can reach hundreds of millions of
rows~\citep{covington2016youtube}. In scientific simulation, it is the
inner loop of iterative solvers for finite-element and finite-difference
discretisations. The same pattern appears in transformer attention with
sparse masks and in the weight-matrix multiply step of pruned and
compressed neural networks~\citep{han2016eie}. Across all of these
workloads the common denominator is that the useful arithmetic is only
a small fraction of the work that a general-purpose processor actually
performs, with most cycles spent chasing indirect indices, handling row
boundaries, and issuing scattered memory accesses.

With the rapid rise in demand for AI, there has been a matching rise in
demand for better model algorithms, more data, and above all more computation.
On the compute side, most AI workloads are now executed on
application-specific integrated circuits (ASICs), since these have been shown to
maximise throughput, energy efficiency, and cost-effectiveness for the
particular matrix operations that AI models need to run. Some of the most
well-known commercial examples are the Tensor Processing Unit (TPU) from
Google and the Neural Processing Unit (NPU) from Apple. The design of any
such ASIC, however, almost always begins with prototyping on a Field
Programmable Gate Array (FPGA), because FPGAs can execute many operations
in parallel and are easily reconfigurable, which keeps architectural
exploration inexpensive and quick. An
FPGA-hosted accelerator for SpMM is therefore a natural first step towards
a fully custom silicon implementation, and is the platform targeted by
this project.

\section{Project Aim}

The aim of this project is to design, implement, and evaluate a
RISC-V-compatible hardware accelerator for compressed sparse row (CSR) SpMM on
the Terasic DE1-SoC FPGA board. The accelerator must integrate cleanly
into a small system-on-chip (SoC), share on-chip block RAM with the CPU,
and deliver a measurable, reproducible speedup over a comparable software
implementation while remaining understandable at the register-transfer
level (RTL) and transparent in its benchmarking methodology. A secondary
aim is to isolate the individual contribution of each optimisation step,
so that the dominant performance bottleneck can be identified through
direct measurement rather than assumption. To achieve this, the design
is developed in three stages (a baseline accelerator, a parallel
DSP-oriented multiply-accumulate enhancement, and a symmetric INT8
packing stage), forming a controlled experimental path in which only
one design dimension is changed at a time.

Beyond the core aim, the project defined three stretch goals at the
outset. The first was to introduce a custom RISC-V instruction to trigger
the accelerator directly from the CPU pipeline, rather than through a
memory-mapped register write, which would reduce launch overhead and
demonstrate that the PicoRV32 soft core can be extended at the instruction
set level. The second was to implement runtime INT8 quantisation so that
operand fetch bandwidth could be reduced by packing four quantised values
into a single 32-bit RAM word. The third was to extend the accelerator
with optional activation functions and a max-pooling stage, so that it
could serve as a building block for a small inference pipeline rather
than as a standalone kernel. Of the three, the INT8 quantisation goal
was achieved and is reported in full in this work; the remaining two
are discussed as future work in Chapter~\ref{ch:conclusion}.

\section{Report Structure}

The remainder of this report is structured as follows.
Chapter~\ref{ch:motivation} reviews the background needed to justify the
design choices, covering SpMM workloads, sparse data formats, prior
accelerator work and the rationale for FPGAs.
Chapter~\ref{ch:platform} defines the implemented SoC platform and its
memory-mapped I/O interface. Chapter~\ref{ch:design} presents the
accelerator architecture and each optimisation step in detail.
Chapter~\ref{ch:methodology} defines the hardware verification approach
and the unit and integration testbenches used to check the RTL.
Chapter~\ref{ch:firmware} describes the bare-metal firmware, the
benchmark suite, and the measured cycle-count and correctness results
across the three optimisation stages. Chapter~\ref{ch:reflection}
discusses lessons learned and project management, and
Chapter~\ref{ch:conclusion} closes the report with a summary and proposed
future work.
